Er verschenen lichten.
Ze gaven inzichten
en goede berichten

aan Emmy deze keer.
Ze voelde wel afkeer
voor het dreigend onweer,

dat spoedig op komst was.
Ze versnelde haar pas.
Toch tierde de judas,

“Daar is mijn eigen stad.
Ik lust stilaan wel wat.
Dus beweeg je neukgat!”

Lichten werden ramen,
toen ze rennend samen
bij de poort aankwamen.

De stadspoort was vrij breed
en met beelden bekleed.
Er was vaardig gesmeed,

in staal als medium,
“Ave imperium!”
en “Pandemonium”

als de welkomstwoorden.
Wat velen niet hoorden,
waren zij die stoorden.

Emmy hoorde ruisen,
via stalen buizen
uit hun houten huizen.

De woeste demonen,
die hels kwamen honen
door haar te bewonen.

Hier woonde de sekte,
die moeder afbekte
en afdeed als gekte.

“Ik wil dat je geknield
door mijn godsstad freewheelt!",
zei de jager bezield.

Deze vernedering
was een verademing
tegen betutteling.

De druk volmaakt te zijn.
Publiekelijke pijn
door schande bleek ook fijn.

Het was zelfs erotisch.
Dat was weer ironisch,
omdat ze platonisch

als non wilde leven.
Sex wilde opgeven
als kunstmatig streven.

Het schandelijke bleek
voor de liefdesdaadleek
een nieuwbakken insteek.

Haar staat was onheilig,
toch voelde het veilig,
zelfs enigszins geilig.

Emmy was dus beslist
een - wat ze nog niet wist -
“exhibitionist”.

Een stadspijpster kwam blij
met haar schelle schalmei
bij het tweetal langszij.

Haar houten herdersfluit
met zijn scherpe geluid
verkondigde zo luid

de jager zijn thuiskeer.
Men was nog in de weer
in de Sint Maarten sfeer.

Dus mensen dromden snel samen rondom de del,
kruipend in haar schaapsvel.

Nog een stadsmuzikant
kwam langs Emmy’s zijkant,
zijn muziek penetrant

in haar penetreren.
Emmy’s degraderen
kon men zeer waarderen.

Iemand stal de schaapsvacht,
die ze losjes meebracht.
Men zag haar volle pracht.

Zonder op te tuigen
vulden de zintuigen,
als stille getuigen

van Emmy haar knapheid,
met ingenomenheid.
Ze kiemden geiligheid.

Gouden triomfbogen,
net boven haar ogen,
werden recht gebogen.

Tonend haar bedenking
over de clustering
en de openbaring.

Men omsingelde haar.
De blijdschap was hoorbaar.
Het was goed nieuws blijkbaar.

Het wachten op een vangst
duurde ditmaal het langst.
Men negeerde de angst

in haar ogen lichtbruin.
De haren op haar kruin
hingen vervuild en schuin

over Emmy’s borstjes,
geplakt door de korstjes
van eerdere morsjes.

Ze was voor hen een ding.
Wat er in haar omging,
kreeg geen belangstelling.

In mistige sferen
met haar paraderen,
zonder vacht of kleren,

gaf de jager hoogmoed.
“Dat doet mij vast geen goed”,
dacht Emmy’s zwaar gemoed.

Op klinkers en geboeid,
bevroren en vermoeid
en net nog flink gestoeid,

ging het kruipen niet snel.
Soms werd Emmy onwel
van haar slopende hel.

De jager dacht terug
aan haar eerste kutplug,
toen Emmy op haar rug

het op hem had gemunt
en typeerde als rund.
Hij bedacht zo een stunt.

Hij vroeg aan een veeboer,
waarlangs haar kruipweg voer,
of hij voor zijn staathoer

kort een koe mocht lenen.
Hij spaarde haar benen,
die afgetrapt schenen.

Voor het vervolgreisje
tilde hij zijn prijsje,
het ontmaagde meisje,

zo witjes en zo moe
boven op de melkkoe.
De menigte keek toe.

Emmy zat verloren
en achterstevoren.
Emmy greep snel een hoorn

om niet af te glijden
en een val te mijden,
lopende het rijden.

Bij de horens pakken,
deed haar lijf niet plakken.
Haar onderuit zakken,

kantelde haar kutje,
met daarop zijn prutje.
Men had op het trutje

volledige inkijk.
Men haalde toen beeldrijk
Emmy sluw door het slijk.

De toeschouwers rondom
lachten zich allen krom
om het gedroogde hom.

Die vlekken lieten zien,
dat hij de rode trien
zeker niet had ontzien.

Emmy was, naast jachtbuit,
ook zijn geschaakte bruid.
Het lokte schelden uit

en een drang tot tuchten.
Emmy moest toen zuchten
onder rotte vruchten.

De jongste stadsliedjes
zongen als bandietjes
hun Sint Maarten liedjes

tijdens hun rondvaarten.
“Sint Maarten, Sint Maarten,
koeien hebben staarten,

meisjes hebben voortaan
niet meer hun rokjes aan.
Daar zit Martinus aan!”

Keer op keer, ongepast,
werd ze intiem betast.
Haar lichaam werd verrast

met steeds een nieuwe hand.
De jager aangebrand
duwde ze aan de kant.

Emmy was zijn bezit.
Daarom hield hij haar klit
stevig in het gelid.

De jager met kabaal
vertelde zijn verhaal.
Men volgde hen massaal.

Sommige meelopers
werden al spekkopers
als haar lichaamstropers.

Want een jongeling mocht
niet mee met de optocht.
Waardoor men een weg zocht,

gevangen door haar blik,
hulpeloos in een strik,
voor hun kloppende pik.

Ze trokken zich dus af.
Het witsel uit hun staf
smeerden ze daarna laf

op het sippe gezicht
van het wankele wicht.
Mentaal door dit gezwicht,

deed het Emmy wel wat,
haar tronie zo beklad
met het jongetjes nat.

Ze wou anders hopen.
In de kijker lopen
en haar onschuld slopen,

was wel lekker misschien.
Ze voelde zich gezien.
Gewaardeerd bovendien.

Al die jonge knullen
met hun stijve lullen
konden haar onthullen

dat ze bekoring bracht
in deze vroege nacht.
Ze hielden van haar pracht.

Net als de mensjagers
en ooit haar seksplagers.
De jonge belagers,

die wekelijks spoten
als ze weer genoten
van haar te ontbloten.

Lopende dit stoetje
landde meer witgoedje
op haar spikkelsnoetje.

Ze voelde zich beschermd,
want ze werd ongeremd
er tegen afgeschermd.

De jager hield tegen
de vluchtige vegen.
Ze konden het plegen

door een kat en muisspel.
Ze waren veel te snel,
maar dat beviel haar wel.

Ook de meisjes vooraan.
Die moedigden hen aan
om vooral door te gaan

met hun vuige makkes.
“Duw eens al dat jakkes
in haar paffe bakkes.”

Dat deed de jeugd prima.
Haar gelaat was daarna
vol met jongenssperma.
